\documentclass[tikz,border=6pt]{standalone}
\usepackage{ctex}
\usepackage{amsmath}
\usepackage{pgfplots}
\pgfplotsset{compat=1.18}

\begin{document}
\begin{tikzpicture}

\begin{axis}[
  width=11cm, height=8cm,
  xmin=0.5, xmax=50,
  ymin=0.0003, ymax=20,
  xlabel={迭代次数 $k$},
  ylabel={$f(\mathbf{x}_k) - f^*$（目标间隙）},
  xlabel style={font=\small},
  ylabel style={font=\small},
  xmode=normal,
  ymode=log,
  axis lines=left,
  axis line style={-Stealth},
  tick style={font=\small},
  xtick={0,10,20,30,40,50},
  grid=both,
  grid style={gray!15},
  legend pos=south west,
  legend style={font=\small, cells={anchor=west}, fill=white, draw=gray!50},
  clip=false,
]

% ISTA: O(1/k) → f(x_k) - f* = C1/k，取 C1=10
\addplot[
  blue, thick,
  domain=1:50, samples=50,
] {10/x};
\addlegendentry{ISTA：$O(1/k)$}

% FISTA: O(1/k^2) → f(x_k) - f* = C2/k^2，取 C2=10
\addplot[
  red!80!black, thick,
  domain=1:50, samples=50,
] {10/(x*x)};
\addlegendentry{FISTA：$O(1/k^2)$}

% ── 参考渐近线（灰色虚线）────────────────────────────────────────────────────
% 1/k 渐近线（斜率 -1 in log-log，但我们用 ymode=log 和 xmode=normal）
% 改用注释箭头标注斜率

% ── 关键点标注 ────────────────────────────────────────────────────────────────
% k=10 时的差距
\fill[blue!80] (axis cs:10, 1.0) circle (2.5pt);
\fill[red!80!black] (axis cs:10, 0.1) circle (2.5pt);
\draw[gray!60, dashed] (axis cs:10, 0.0003) -- (axis cs:10, 1.2);
\node[font=\small, anchor=south, align=center] at (axis cs:10, 1.25)
  {$k=10$};
\draw[{Stealth[length=4pt]}-{Stealth[length=4pt]}, gray!70]
  (axis cs:11.5, 0.1) -- (axis cs:11.5, 1.0);
\node[font=\small, gray!60, anchor=west] at (axis cs:11.8, 0.3)
  {10倍差距};

% k=30 时的差距
\fill[blue!80] (axis cs:30, 0.333) circle (2.5pt);
\fill[red!80!black] (axis cs:30, 0.0111) circle (2.5pt);

% ── 斜率标注 ─────────────────────────────────────────────────────────────────
% ISTA 斜率标注（log-y 中 -1/k 是非线性，直接注释）
\node[font=\small, blue, anchor=south west] at (axis cs:20, 0.55)
  {$\sim \dfrac{1}{k}$};

\node[font=\small, red!80!black, anchor=north east] at (axis cs:22, 0.025)
  {$\sim \dfrac{1}{k^2}$};

% ── 收敛速度比较说明 ────────────────────────────────────────────────────────
\node[font=\footnotesize, align=left, draw=gray!40, fill=white!90,
      rounded corners, inner sep=4pt]
  at (axis cs:36, 0.5) {
    \textbf{每步计算量相同}\\
    FISTA 仅增加\\
    $\mathbf{y}_k,\,t_k$ 两变量
  };

% ── 精度要求线（ε = 0.01）─────────────────────────────────────────────────────
\addplot[
  gray!50, dashed, thick,
  domain=0.5:50,
] {0.01};
\node[font=\small, gray!60, anchor=west] at (axis cs:42, 0.012)
  {精度 $\varepsilon$};

% ISTA 到达 ε 时的迭代数（10/k = 0.01 → k=1000）：超出范围，只标注近似
\draw[-{Stealth[length=4pt]}, gray!50, dashed]
  (axis cs:50, 0.2) -- (axis cs:50, 0.01);
\node[font=\tiny, gray!60, anchor=north] at (axis cs:50, 0.009)
  {ISTA 需 $\sim\!1000$ 步};

% FISTA 到达 ε 时（10/k^2 = 0.01 → k=32）
\fill[red!80!black] (axis cs:31.6, 0.01) circle (2pt);
\node[font=\small, red!80!black, anchor=south]
  at (axis cs:31.6, 0.012) {$k\approx 32$};

% 图标题
\node[font=\small, align=center] at (axis cs:25, 7)
  {ISTA vs FISTA 收敛速度对比（对数纵轴）};

\end{axis}
\end{tikzpicture}
\end{document}
